%!TEX root = ../LaBoDither_springer_v0.tex
%!TEX encoding = UTF-8
%!TEX spellcheck = en_US

%%%% ---- LaTeX File ---- %%%%
\providecommand{\PathP}{../../Parts}  % Path to Parts folders
\providecommand{\PathR}{../../Resources}  % Path to  Resources folders
\providecommand{\PathS}{../../Supplies/}  % Path to Supplies folders
\providecommand{\PathSO}{../../Sources/}  % Path to Sources folders
\providecommand{\lang}{english}  % english, french
\providecommand{\BuildTrigger}{buildnew}  % Trigger to build or not standalone parts.

\makeatletter
	\def\input@path{ {\PathSO}{./another-directory/} }  % to be able to find package in source folders.
		\ifx\c@levelstanda\undefined  % to create a newcounter if it doesn't exists.
		\newcounter{levelstanda}
	\else
		{}% it already exist
	\fi
\makeatother

%%%%---- DOCUMENT-CLASS -------------------------%%%%
\documentclass{article}
	\usepackage[l2tabu, orthodox]{nag}  % Helps to correct the script

%%%%---- OS & FONT -------------------------%%%%
	\usepackage[utf8]{inputenc}  % Input keyboard encoding: latin1=windows, utf8=linux, applemac=Mac
	\usepackage[T1]{fontenc}  % Output font encoding: T1 for Europeans languages.


%%%%---- PERSONAL Custom ----------------------%%%%
		\newcommand{\Pathbasebib}{\PathR/ref_RN.bib,\PathR/ref_XX.bib}% Careful no space!
		\newcommand{\PathSymbolsEntries}{\PathR/UsersSymbolsAcronyms.bib}
	\input{\PathR/UsersPackagesDefs.tex}
	% \input{\PathR/UsersSymbolsAcronyms.tex}
	\input{\PathR/UsersMathsCommands.tex}

%%%%---- ONLY to ADAPT without JOURNAL CLASS ----%%%%
	\usepackage[margin=3cm]{geometry}
	\usepackage{natbib}
		\bibliographystyle{plain}


%%%%---- Journal Spec -------------------------%%%%
\title{Cover Letter}
\author{
	Romain~\famName{Noël}, John~\famName{Doe}
}


%%%%%%%%%%%%%%%  MAIN DOCUMENT  %%%%%%%%%%%%%%%%%%%%%
\begin{document}
\stepcounter{levelstanda}

	\begin{flushleft}
		\underline{Origin}:\newline
		Romain~\famName{Noël} \newline
		Gustave Eiffel University \newline
		romain.noel@univ-eiffel.fr \newline
		10th of January 2024
	\end{flushleft}

	\begin{flushright}
		\underline{Destination}: \\
		Pr. Ahmet Enis Cetin, Editor in Chief \\
		SIVP Journal
	\end{flushright}

	\begin{center}
		\underline{Subject}: Submission of the article entitled: "A Generalized Dithering Approach using the Lattice Boltzmann Method"
	\end{center}


	Dear Editor,

	I am writing to submit our manuscript entitled "A Generalized Dithering Approach using the Lattice Boltzmann Method" for consideration for publication in {Signal, Image and Video Processing}.
	I believe that the research presented in this paper makes a valuable contribution to the field of {signal and image processing}.

	We propose a comprehensive dithering approach employing the Lattice Boltzmann Method (LBM).
	This generalized approach facilitates the exploration of various models, enabling adaptation to specific features and challenges.
	As exemplifications, we introduce and assess four models, encompassing anisotropic gradient-based relaxation time and an improved equilibrium distribution function.
	The primary objective is to enhance the quality and reduce artifacts.
	Implementation of these models on Graphics Processing Units (GPUs) and their generic integration using the OBLiX framework are discussed.
	The outcomes of multiple test cases demonstrate that such models yield superior qualitative and quantitative results, underlining the adaptability of the proposed generalization.
	This paper establishes a framework for dithering using the Lattice Boltzmann Method, allowing algorithm adaptation to intricate scenarios beyond conventional cases typically encountered.

	I believe that the innovative approach presented in this paper aligns well with the scope and interests of your journal.
	The utilization of the Lattice Boltzmann Method for dithering, along with the introduction of novel models, contributes to the understanding of both LBM and dithering methods.
	The experiments conducted, particularly those involving GPU implementation and interfacing with the OBLiX framework, highlight the practical applicability of the proposed approach.

	I think %have ensured
	that the manuscript adheres to the submission guidelines of {Signal, Image and Video Processing}.
	I am confident that the findings presented in this paper will be of interest to the journal's readership and will contribute positively to the ongoing discourse in our field.

	Thank you for considering my submission.
	I look forward to the opportunity for peer review and feedback from the esteemed reviewers at {Signal, Image and Video Processing}.

	\begin{flushright}
		Best regards, \\
		Dr. Romain \famName{Noel} (in the name of all authors), \\
		Scientific research, \\
		Gustave Eiffel University, \\
		romain.noel@univ-eiffel.fr
	\end{flushright}

\ifnum\value{levelstanda}>1
	{}% Not in the main.
\else

	% References
	% \bibliography{\Pathbasebib} % for bibtex or natbib
	% %\printbibliography     % for biblatex

	%%% To comment for the final version
	% %\tableofcontents
	% \printglossary[type=\glsxtrabbrvtype]
	% \printglossary[type=symbol]
\fi

\addtocounter{levelstanda}{-1}
\end{document}